\documentclass[conference]{IEEEtran}
\IEEEoverridecommandlockouts
\usepackage{cite}
\usepackage{amsmath,amssymb,amsfonts}
\usepackage{algorithmic}
\usepackage{graphicx}
\usepackage{textcomp}
\usepackage{xcolor}
\usepackage{url}
\usepackage{booktabs}
\usepackage{listings}

\def\BibTeX{{\rm B\kern-.05em{\sc i\kern-.025em b}\kern-.08em
    T\kern-.1667em\lower.7ex\hbox{E}\kern-.125emX}}

\lstset{
    basicstyle=\fontsize{7}{8}\ttfamily,
    breaklines=true,
    frame=single,
    backgroundcolor=\color{gray!10}
}

\begin{document}

\title{Targeted Bus-Cycle Interception: A Low-Latency WebAssembly Bridge for 16-Bit Bare-Metal Audio Synthesis and Multi-Timbral Acoustic Resynthesis}

\author{
\IEEEauthorblockN{Krishna Aher\IEEEauthorrefmark{1}, Sanskar Bhargude\IEEEauthorrefmark{1}, Ghansham Agaldare\IEEEauthorrefmark{1}, Hari Birare\IEEEauthorrefmark{1}, Akash Kumar\IEEEauthorrefmark{1}, and Prof. Gopal Upadhye\IEEEauthorrefmark{2}}
\IEEEauthorblockA{\IEEEauthorrefmark{1}Department of Multidisciplinary Engineering / Artificial Intelligence \& Data Science\\
\IEEEauthorrefmark{2}Faculty Project Guide, Department of Multidisciplinary Engineering\\
Vishwakarma Institute of Technology (Autonomous Institute Affiliated to Savitribai Phule Pune University), Pune, Maharashtra, India}
}

\maketitle

\begin{abstract}
The 16-bit 8086 microprocessor architecture remains an essential pedagogical foundation in computer engineering curricula, embedded systems design, and deterministic hardware-software co-design. However, legacy x86 audio execution is historically constrained by two divergent extremes: physical bare-metal execution is restricted to harsh monophonic 1-bit square waves, whereas existing browser-based emulators (such as v86 and DOSBox-Wasm) impose a severe ``Virtualization Tax'' requiring over 140~MB of memory and introducing 60 to 120~ms of buffer dispatch jitter. In this paper, we present \textit{Crimson Orbit}, a novel dual-state computing platform that bridges bare-metal real-mode 8086 assembly with modern WebAudio Digital Signal Processing (DSP). Rather than simulating an entire motherboard, our approach introduces \textbf{Targeted Bus-Cycle Interception}: a lightweight WebAssembly execution bridge that traps I/O port writes exclusively to the Intel 8253 Programmable Interval Timer (Port 42h) and Intel 8255 Programmable Peripheral Interface (Port 61h), serializing them into an ultra-compact 4-byte micro-packet ($\langle$CMD, DIV\_LO, DIV\_HI, DURATION$\rangle$). We implement an interactive 60-note in-RAM circular tape sequencer (\texttt{composer.asm}) executing bare-metal without operating system interrupts, alongside a bifurcated dual-engine synthesis pipeline supporting cycle-accurate 1-bit Galois LFSR square waves and 44.1~kHz multi-timbral acoustic resynthesis (Steinway Grand Piano, Martin Acoustic Guitar, Ludwig Studio Drums). Empirical benchmarking across 1,000 iterations confirms that Crimson Orbit achieves a mean bus-trapping latency of 0.889~$\mu$s, an end-to-end dispatch latency of 11.8~ms (a 6.3$\times$ improvement over full virtual machines), and a runtime memory footprint of just 12.4~MB (an 82\% to 91\% reduction).
\end{abstract}

\begin{IEEEkeywords}
WebAssembly, 8086 Real-Mode Assembly, Bus-Cycle Interception, Low-Latency Audio, Web Audio API, Microprocessor Virtualization, Acoustic Resynthesis, Computer Systems Pedagogy.
\end{IEEEkeywords}

\section{Introduction \& Historical Motivation}

\subsection{The 1981 IBM PC Audio Subsystem}
Undergraduate computer engineering curricula mandate hands-on experiential mastery of low-level machine organization, peripheral bus input/output (I/O) interfacing, and cycle-accurate execution, as formalized in the ACM/IEEE Computer Science Curricula 2023 (CS2023) guidelines~\cite{cs2023}. For over four decades, the Intel 8086 real-mode architecture has remained the premier pedagogical vehicle because of its architectural transparency, deterministic register set, and absence of preemptive kernel abstractions.

However, bare-metal audio generation on historical x86 hardware embodies an extreme physical bottleneck. When IBM launched the Personal Computer (Model 5150) in August 1981, audio capabilities were intentionally minimalist. The platform coupled Channel 2 of the Intel 8253 Programmable Interval Timer (PIT) and Bit 0/1 of Port 61h on the Intel 8255 Programmable Peripheral Interface (PPI) directly to a small, unamplified 2-inch permanent-magnet cone~\cite{collins2020}.

Because the PPI drives the speaker voice coil using binary transistor-transistor logic (TTL) switching (+5V logic high vs. 0V logic low), physical bare-metal x86 software cannot produce arbitrary polyphony or continuous waveforms without complex pulse-width modulation hacks that overwhelm the 4.77 MHz CPU. Instead, timer-driven tone generation produces pure square waves governed by the Fourier series:
\begin{equation}
x(t) = \frac{4}{\pi} \sum_{k=1}^{\infty} \frac{1}{2k - 1} \sin((2k - 1)\omega_0 t)
\end{equation}
This distribution contains infinite odd harmonics whose amplitudes decrease as $1/n$, yielding an acoustic output with a Total Harmonic Distortion (THD) of 48.34\%. Consequently, bare-metal assembly music is universally characterized by harsh, piercing ``chiptune beeps.''

\subsection{Modern Silicon Dilemma \& The Virtualization Tax}
In contemporary engineering classrooms, physical 8086 machines and ISA bus expansion slots have vanished. To provide access to vintage x86 environments, educators have adopted two primary alternatives: software emulators and full-system browser virtual machines.

Standalone emulators such as EMU8086~\cite{emu8086} provide register-level debugging but suffer from a severe pedagogical disconnect: students watch static hexadecimal values increment in \texttt{AX} and \texttt{DX} without intuitive sensory feedback connecting code execution to physical peripheral behavior.

Conversely, modern browser-based virtualization systems such as \textit{v86}~\cite{v86} and \textit{DOSBox-Wasm}~\cite{dosbox} compile legacy C/C++ x86 emulators to WebAssembly. While functionally capable of booting entire operating systems, these monolithic engines incur a heavy \textbf{``Virtualization Tax''}. They allocate 64 to 150 MB of contiguous linear memory to simulate floppy disk controllers (Intel 8272A), DMA controllers (Intel 8237), VGA display framebuffers, and protected-mode memory management units. Crucially, their audio generation routines are tied to the 60 Hz display refresh loop, introducing 60 to 120 ms of buffer dispatch jitter. Because human psychoacoustics mandates an end-to-end latency below 15 to 20 ms for real-time interactive musical response~\cite{wilson2002}, full virtual machines remain unsuitable for responsive audio synthesis.

\subsection{Problem Statement \& Research Objectives}
The central research problem addressed in this work is: \textit{How can we bridge cycle-accurate 16-bit real-mode x86 assembly execution with low-latency, multi-timbral digital audio synthesis in modern distributed browser runtimes without incurring the memory bloat and latency penalties of monolithic hardware virtualization?}

To resolve this problem, we establish four concrete engineering objectives:
\begin{enumerate}
\item \textbf{Sub-12ms End-to-End Latency:} Eliminate display-frame audio quantization by trapping bus cycles at the machine-cycle level and dispatching to a dedicated \texttt{AudioWorklet} thread.
\item \textbf{Overhead Reduction $>$80\%:} Replace full motherboard emulation with targeted peripheral trapping, slashing RAM overhead below 15 MB.
\item \textbf{Acoustic Transduction:} Provide a mathematically rigorous pipeline mapping discrete timer countdowns to high-fidelity multi-timbral PCM instruments.
\item \textbf{Bare-Metal In-RAM DAW:} Implement an autonomous 60-note circular composition sequencer in assembly without DOS \texttt{INT 21h} interrupts.
\end{enumerate}

\begin{figure}[t]
\centering
\includegraphics[width=\columnwidth]{Assets/Images/CrimsonOrbit_System_Architecture.jpg}
\caption{Three-tier system architecture of Crimson Orbit: Layer 1 Bare-Metal 8086 Engine, Layer 2 Low-Latency WebAssembly Bus Interceptor Bridge, and Layer 3 Multi-Timbral Acoustic Resynthesis Web Studio.}
\label{fig:arch}
\end{figure}

\section{System Architecture \& 8086 Kernel Internals}

\subsection{Autonomous Real-Mode Memory Layout}
To achieve total execution autonomy without host operating system dependencies, Crimson Orbit operates as a bare-metal flat binary booted via a custom 512-byte Master Boot Record (\texttt{boot.asm}). The CPU memory segmentation model establishes strict deterministic partitions:
\begin{itemize}
\item \texttt{0000:7C00h} -- \texttt{0000:7DFFh}: MBR Bootloader (512 Bytes)
\item \texttt{1000:0000h} -- \texttt{1000:3ABFh}: Executable Flat Kernel Image (15 KB)
\item \texttt{1000:2000h} -- \texttt{1000:23FFh}: In-RAM Tape Sequencer Buffer (1 KB)
\item \texttt{9000:0000h} -- \texttt{9000:FFFFh}: System Stack Segment (64 KB)
\end{itemize}

\subsection{Hardware Register Programming}
The audio generation hardware subsystem is governed by two integrated peripheral controller chips: the Intel 8253 PIT and the Intel 8255 PPI.

The Intel 8253 PIT incorporates three independent 16-bit countdown counters clocked by a master oscillator derived from the color burst crystal:
\begin{equation}
f_{\text{osc}} = \frac{14.31818\text{ MHz}}{12} = 1,193,181.81\text{ Hz} \approx 1,193,180\text{ Hz}
\end{equation}

Channel 2 of the PIT is dedicated to sound synthesis. To initialize Channel 2, the CPU writes control word \texttt{0B6h} (binary \texttt{10110110b}) to Control Port 43h:
\begin{itemize}
\item Bits 7--6 = \texttt{10b}: Select Counter 2
\item Bits 5--4 = \texttt{11b}: Read/Load LSB followed by MSB
\item Bits 3--1 = \texttt{011b}: Mode 3 (Square Wave Generator)
\item Bit 0 = \texttt{0b}: 16-bit Binary Counter
\end{itemize}

To synthesize a musical frequency $f$, the 16-bit divisor is computed via:
\begin{equation}
\text{Divisor} = \lfloor 1,193,180 / f + 0.5 \rfloor
\end{equation}

The divisor is transferred to PIT Data Port 42h in two consecutive byte writes: Least Significant Byte (LSB) followed by Most Significant Byte (MSB).

Gating of the physical speaker cone is governed by System Control Port B (Port 61h) on the Intel 8255 PPI. Bit 0 controls GATE 2 (enabling the counter), while Bit 1 gates the counter output to the speaker cone. Enabling sound requires setting both bits high without disturbing existing system status bits:

\begin{lstlisting}[language={[x86masm]Assembler}]
; Intel 8253/8255 Tone Driver (Source/speaker.asm)
mov dx, 0012h ; High word of 1,193,180
mov ax, 34DCh ; Low word of 1,193,180
div bx        ; AX = Divisor (1,193,180 / BX_Freq)
mov al, 0B6h  ; Channel 2, Mode 3, LSB/MSB
out 43h, al   ; PIT Command Register
mov al, bl
out 42h, al   ; Send Divisor LSB
mov al, bh
out 42h, al   ; Send Divisor MSB
in  al, 61h   ; Read PPI Port B
or  al, 03h   ; Assert Bits 0 & 1 (Gate & Speaker)
out 61h, al   ; Enable Sound Output
\end{lstlisting}

\subsection{Galois LFSR Pseudo-Random Noise Synthesis}
To synthesize percussive transients (snare hits, cymbals, hi-hats) on bare-metal hardware without digital-to-analog converters, \texttt{speaker.asm} incorporates a 16-bit Galois Linear Feedback Shift Register (LFSR). The LFSR implements the maximal-length polynomial:
\begin{equation}
P(x) = x^{16} + x^{14} + x^{13} + x^{11} + 1
\end{equation}
With feedback tap mask \texttt{0ACE1h} (or \texttt{0B400h}), the generator executes an in-register shift loop. When the least significant bit is set, the accumulator is XORed with the tap mask, and Bit 1 of Port 61h is rapidly toggled. This produces authentic band-limited white noise directly from bare-metal assembly.

\section{Targeted Bus-Cycle Interception Bridge}

\subsection{Micro-Interception Theory}
Rather than executing an entire virtualized motherboard with memory paging, DMA engines, and display refresh scans, Crimson Orbit introduces \textbf{Targeted Bus-Cycle Interception}. The execution bridge instruments the 8086 CPU execution loop exclusively at the point of I/O port instruction decoding (opcodes \texttt{0xE6}, \texttt{0xE7}, \texttt{0xEE}, \texttt{0xEF}).

When the CPU executes an instruction targeting Port 42h (PIT Channel 2) or Port 61h (PPI Speaker Gate), the bus trap triggers. The trap intercepts the register state, extracts the divisor and duration values, and serializes the hardware bus transaction into an immutable 4-byte micro-packet.

\subsection{The 4-Byte Micro-Packet Protocol}
The micro-packet specification is engineered for minimal serialization overhead and zero memory allocation. Each intercepted event is packed into exactly 32 bits:
\begin{equation}
P = \langle \text{CMD}, \text{DIV\_LO}, \text{DIV\_HI}, \text{DURATION} \rangle
\end{equation}
\begin{itemize}
\item \textbf{Byte 0 (CMD, 8 bits):} Command opcode: \texttt{0x00} = NOTE\_OFF (Port 61h bits cleared); \texttt{0x01} = NOTE\_ON (Tone active); \texttt{0x02} = NOISE\_BURST (Galois LFSR percussion); \texttt{0x03} = TEMPO\_SYNC.
\item \textbf{Byte 1 (DIV\_LO, 8 bits):} Divisor Least Significant Byte written to Port 42h.
\item \textbf{Byte 2 (DIV\_HI, 8 bits):} Divisor Most Significant Byte written to Port 42h.
\item \textbf{Byte 3 (DURATION, 8 bits):} Step duration encoded in 10 ms discrete quantization intervals ($10\text{ ms} \le \Delta t \le 2550\text{ ms}$).
\end{itemize}

The browser engine deserializes the packet instantaneously:
\begin{equation}
\text{Divisor} = (\text{DIV\_HI} \ll 8) \mid \text{DIV\_LO}
\end{equation}
\begin{equation}
f_{\text{synth}} = \frac{1,193,180}{\text{Divisor}}\text{ Hz}
\end{equation}

\subsection{Zero-Copy SharedArrayBuffer \& AudioWorklet}
To ensure deterministic sub-12ms delivery without JavaScript garbage collection interruptions, micro-packets are posted to a lock-free circular ring buffer backed by a \texttt{SharedArrayBuffer}. Synchronization between the Wasm execution thread and the dedicated WebAudio real-time audio thread is managed via lock-free atomic operations (\texttt{Atomics.wait} and \texttt{Atomics.notify}). The audio thread consumes packets in 128-sample quanta (2.90 ms frames), ensuring smooth, jitter-free playback.

\section{Multi-Timbral Acoustic Resynthesis Pipeline}

\subsection{Control-Plane vs. Synthesis-Plane Decomposition}
A foundational theoretical insight of Crimson Orbit is the formal separation of computer sound into two distinct planes:
\begin{enumerate}
\item \textbf{Control Plane (Bare-Metal 8086 Assembly):} Executes sequencing logic, note duration timers, pitch divisor arithmetic, and in-RAM composition storage.
\item \textbf{Synthesis Plane (WebAudio DSP):} Operates as the high-resolution acoustic transducer and digital-to-analog converter (DAC).
\end{enumerate}

Historically, the 8086 processor never synthesized analog sound waves directly; it commanded peripheral timers and dedicated sound chips (e.g., Creative Labs Sound Blaster, AdLib). Crimson Orbit faithfully preserves the 8086 assembly control plane while modernizing the acoustic synthesis plane.

\begin{figure}[t]
\centering
\includegraphics[width=\columnwidth]{Assets/Images/fig3_fft_harmonic_spectrum.png}
\caption{Time and frequency domain characterization: 1-bit PIT square wave exhibiting odd harmonics and THD = 48.3\% versus acoustic resynthesis exhibiting natural formant decay and THD $<$ 1.2\%.}
\label{fig:fft}
\end{figure}

\subsection{Dual-Engine Synthesis Mode}
To satisfy both historical purists and modern musicians, Crimson Orbit incorporates a hardware engine switch:

\textit{1) Authentic 1-Bit Mode (Zero Audio Samples):} In this mode, zero external audio samples or soundfonts are loaded. The browser creates a pure WebAudio \texttt{OscillatorNode} generating mathematical square waves derived directly from the intercepted divisor. Percussion is synthesized on-the-fly using a software implementation of the 16-bit Galois LFSR with seed \texttt{0ACE1h}, replicating the exact acoustic buzz of an IBM 5150 PC speaker.

\textit{2) Acoustic Resynthesis Mode:} The intercepted frequency and duration parameters are routed to a multi-timbral PCM acoustic resynthesis engine. The engine models three primary concert acoustic instruments:
\begin{itemize}
\item \textbf{Steinway Concert Grand Piano:} Models strike hammer transients, string decay envelopes, and harmonic wooden soundboard resonances.
\item \textbf{Martin D-28 Acoustic Dreadnought Guitar:} Captures polyphonic string plucking, fretboard damping, and acoustic cavity formants across standard guitar tunings (E2 to E4).
\item \textbf{Ludwig Super Classic Studio Drum Kit:} Translates Galois LFSR noise bursts into high-energy snare wire snaps, bass drum kicks, and brass crash cymbals.
\end{itemize}

\subsection{Real-Time FFT Spectral Visualizer}
The web studio incorporates a 2048-point Fast Fourier Transform (FFT) real-time spectrum analyzer operating at 60 FPS. As depicted in Fig.~\ref{fig:fft}, the 1-bit square wave exhibits infinite odd harmonics ($f_0, 3f_0, 5f_0, 7f_0, \dots$) confirming a 48.34\% THD, while the acoustic resynthesis model exhibits warm physical resonance formants with THD $<$ 1.2\%.

\section{In-RAM Tape Sequencer (\texttt{composer.asm})}

\subsection{Eliminating Operating System Services}
Conventional DOS music programs rely on DOS software interrupts (\texttt{INT 21h}) to read file arrays from disk. To achieve true bare-metal autonomy, Crimson Orbit implements an interactive 60-note circular memory tape sequencer running in pure real mode (\texttt{Source/composer.asm}).

The composer module allocates a dedicated 1 KB buffer in the data segment (\texttt{Song\_CustomUser}). Each musical step is encoded as an immutable 32-bit word pair:

\begin{lstlisting}[language={[x86masm]Assembler}]
; Circular Tape Memory Buffer (Source/composer.asm)
Song_CustomUser:
    dw 440, 300    ; Step 1: Note A4 (440 Hz), 300 ms
    dw 523, 600    ; Step 2: Note C5 (523 Hz), 600 ms
    dw 659, 300    ; Step 3: Note E5 (659 Hz), 300 ms
    dw 0FFFFh, 0   ; Sentinel: End of Sequence Marker
\end{lstlisting}

\subsection{Bidirectional Memory Deserializer}
The browser studio incorporates a bidirectional binary deserializer. When an assembly program updates the tape buffer in memory, the bridge uses \texttt{DataView.getUint16(offset, true)} to deserialize the word pairs into an interactive 8-track, 16-step matrix sequencer in the DOM. Conversely, edits performed in the web UI are encoded into 16-bit little-endian binary words and written back into the execution image, demonstrating true bidirectional hardware-software synchronization.

\begin{table}[t]
\caption{Empirical Performance Benchmarking (1,000 Iterations)}
\label{tab:benchmarks}
\centering
\begin{tabular}{lcccc}
\toprule
\textbf{Metric} & \textbf{v86 VM} & \textbf{DOSBox} & \textbf{Crimson Orbit} & \textbf{Benefit} \\
\midrule
RAM Footprint & 142.6 MB & 68.2 MB & \textbf{12.4 MB} & 91.3\% Lower \\
Latency (Mean) & 74.2 ms & 88.5 ms & \textbf{11.8 ms} & 6.3$\times$ Faster \\
Jitter ($\sigma$) & $\pm 18.4$ ms & $\pm 22.1$ ms & \textbf{$\pm 0.35$ ms} & Sub-Perceptual \\
Bus Trap Time & N/A & N/A & \textbf{0.889 $\mu$s} & Microsecond \\
Packet Size & Monolithic & Monolithic & \textbf{4 Bytes} & Minimal \\
Cold Boot Time & 3.80 s & 2.40 s & \textbf{0.18 s} & Instant Web \\
Binary Payload & $>$15 MB & $>$8 MB & \textbf{14.7 KB} & 98\% Smaller \\
\bottomrule
\end{tabular}
\end{table}

\section{Empirical Evaluation \& Benchmarking}

\subsection{Experimental Methodology}
To rigorously evaluate system performance, we developed an automated test harness (\texttt{Tools/benchmark\_metrics.py}) executing 1,000 continuous bus-trapping iterations across the standard chromatic scale (C3 to C6, 130.81 Hz to 1046.50 Hz). The benchmarking testbed was executed on an Intel Core i7-12700H system running Windows 11 with the Chromium V8 WebAssembly engine. Timings were captured using nanosecond hardware counters (\texttt{time.perf\_counter\_ns()}) and high-resolution WebAudio clocks (\texttt{AudioContext.currentTime}).

\subsection{Memory Footprint Analysis}
As detailed in Table~\ref{tab:benchmarks} and illustrated in Fig.~\ref{fig:memory}, Crimson Orbit achieves an active runtime memory footprint of just \textbf{12.4 MB}. In comparison, v86 consumes \textbf{142.6 MB} and DOSBox-Wasm consumes \textbf{68.2 MB}. This represents an \textbf{81.8\% to 91.3\% reduction} in RAM overhead.

This dramatic efficiency gain directly validates our architectural thesis: by discarding the simulation of video display memory (VRAM), DMA controllers, and IDE drives, the memory requirement is bounded strictly by the bare-metal kernel image (15 KB) and the WebAudio buffer graph.

\begin{figure}[t]
\centering
\includegraphics[width=\columnwidth]{Assets/Images/fig1_latency_jitter.png}
\caption{End-to-end audio dispatch latency and timing jitter across 1,000 continuous iterations. Crimson Orbit achieves 11.8 ms (comfortably below the 15 ms human real-time threshold) with sub-perceptual $\pm 0.35$ ms jitter.}
\label{fig:latency}
\end{figure}

\begin{figure}[t]
\centering
\includegraphics[width=\columnwidth]{Assets/Images/fig2_memory_payload.png}
\caption{System resource allocation and payload comparison: (a) Active runtime memory overhead (12.4 MB vs 142.6 MB, a 91.3\% reduction). (b) Standalone binary payload footprint (14.7 KB bare-metal kernel).}
\label{fig:memory}
\end{figure}

\subsection{Audio Latency \& Dispatch Jitter}
As plotted in Fig.~\ref{fig:latency}, end-to-end audio dispatch latency is reduced from \textbf{74.2 ms} in v86 and \textbf{88.5 ms} in DOSBox-Wasm down to \textbf{11.8 ms} in Crimson Orbit. This represents a \textbf{6.3$\times$ latency reduction}, successfully crossing the critical 15 ms perceptual barrier required for real-time acoustic instrument playability.

Furthermore, dispatch timing jitter ($\sigma$) drops from $\pm 18.4$ ms in v86 down to \textbf{$\pm 0.35$ ms} in Crimson Orbit. This stability is directly attributable to the zero-allocation 4-byte micro-packet protocol and the lock-free \texttt{SharedArrayBuffer} audio ring buffer, which immunizes the synthesis pipeline against JavaScript Garbage Collection (GC) pauses.

\section{Conclusion \& Future Work}
This paper presented Crimson Orbit, a low-latency computing platform uniting bare-metal 16-bit 8086 assembly execution with modern WebAudio DSP. By introducing Targeted Bus-Cycle Interception, the system bypasses the Virtualization Tax of monolithic emulators, achieving an 11.8 ms end-to-end latency, $\pm 0.35$ ms jitter, and a 12.4 MB RAM footprint. Future work will extend the micro-interception architecture to Yamaha YMF262 (OPL3) FM synthesis registers at Port 388h and collaborative multi-user WebRTC jam sessions.

\section*{Acknowledgment}
The authors express their deepest gratitude to project faculty guide Prof. Gopal Upadhye for his invaluable technical guidance, pedagogical insight, and continuous support throughout the architecture and empirical evaluation of this research. We also thank the Department of Multidisciplinary Engineering and Artificial Intelligence \& Data Science at Vishwakarma Institute of Technology, Pune, for providing laboratory computing resources.

\begin{thebibliography}{00}
\bibitem{cs2023} ACM/IEEE Computer Society, ``Computer Science Curricula 2023 (CS2023): Core Guidelines in Architecture and Systems,'' \emph{IEEE/ACM Joint Curriculum Task Force}, IEEE/ACM Press, pp. 1--142, 2023.
\bibitem{collins2020} M. Collins, ``The 1-Bit Beep: Physical Constraints and Creative Practices of Vintage Computer Sound,'' \emph{IEEE Annals of Computing History}, vol. 42, no. 3, pp. 12--25, 2020.
\bibitem{v86} F. Hemmer, ``v86: x86 Hardware Virtualization in WebAssembly and JavaScript,'' \emph{Open Technical Monograph}, GitHub Repository, 2014--2024.
\bibitem{dosbox} B. Cai, ``DOSBox-Wasm: Emulating Retro x86 Binaries in Modern Web Browsers,'' \emph{Technical Report}, WebAssembly Systems Community, 2021.
\bibitem{wilson2002} P. R. Wilson, ``Operating System Support for Low-Latency Musical Computing,'' \emph{IEEE Trans. Software Engineering}, vol. 28, no. 6, pp. 589--601, 2002.
\bibitem{emu8086} R. Kumar and P. Sharma, ``Comparative Analysis of Microprocessor Simulation Tools in Computer Engineering Pedagogy,'' \emph{Computer Applications in Engineering Education}, vol. 30, no. 4, pp. 1102--1118, 2022.
\bibitem{patterson2020} D. A. Patterson and J. L. Hennessy, \emph{Computer Organization and Design: The Hardware/Software Interface (RISC-V Edition)}, Morgan Kaufmann, 2020.
\bibitem{blem2014} H. A. Blem, J. Menon, and K. Sankaralingam, ``A Detailed Analysis of the Architectural Bottlenecks in Low-Latency Emulation,'' in \emph{Proc. IEEE Int. Symp. High Performance Computer Architecture (HPCA)}, pp. 411--422, 2014.
\bibitem{w3c2021} P. Adenot and C. Wilson, ``Web Audio API: W3C Recommendation,'' \emph{World Wide Web Consortium (W3C)}, Dec. 2021.
\bibitem{intel8086} Intel Corporation, ``8086 16-Bit HMOS Microprocessor User's Manual,'' Order No. 9800722-03, Santa Clara, CA, 1981.
\bibitem{intel8253} Intel Corporation, ``8253/8253-5 Programmable Interval Timer Data Sheet,'' Order No. 231308-004, Santa Clara, CA, 1993.
\bibitem{intel8255} Intel Corporation, ``8255A/8255A-5 Programmable Peripheral Interface Data Sheet,'' Order No. 231310-003, Santa Clara, CA, 1993.
\end{thebibliography}

\end{document}
